\section{Automatic and Interactive Proof}
\label{sec:workflow}

\subsection{A library for B reasoning}

The encoding is accompanied by definitions and proved lemmas for the
represented B operators. For example, the domain of a total function is
its source set; a function's image of a finite set is finite; and the
minimum of a finite nonempty image belongs to that image. Such results
serve both interactive proofs and automatic WD discharge. Since sets
and relations use Mathlib representations, users can also apply ordinary
set extensionality, image laws, arithmetic reasoning, and rewriting.

WD automation is extensible through labelled collections of lemmas.
The implementation provides \texttt{wd\_min}, \texttt{wd\_max},
\texttt{wd\_app}, and \texttt{wd\_card} collections and corresponding
tactics. A theorem registered with a label becomes available to the
associated search. For a linearly ordered codomain, such as the integers
used in the case study, the minimum collection contains a rule of the form
\[
  f\in A\pfun B\quad\land\quad
  S\in\operatorname{FIN}_1(\operatorname{dom}(f))
  \quad\Longrightarrow\quad \WD_{\min}(f[S]),
\]
where $\operatorname{FIN}_1(X)$ denotes the finite nonempty subsets of $X$.
The same pattern supports maximum; cardinality uses finiteness rules;
application uses function properties and domain membership.

The user-facing \texttt{barrel\_solve} tactic combines basic introduction
and triviality steps, B typing tactics, specialized WD reasoning, and
Lean's \texttt{grind} tactic. The specialized tactics introduce assumptions,
substitute equalities, and generalize proof arguments before searching
for an applicable lemma. This reduces sensitivity to the particular
proof term supplied as a WD argument. All successful attempts construct
Lean proofs; a failed attempt leaves a residual obligation for the user.

A domain-specific lemma can connect an invariant to a reusable WD rule.
In the monitoring case study, a helper proves that a permitted zone's
available sensors form a finite nonempty subset of the reading map's domain.
The general finite-image rules then establish the extrema conditions.
Section~\ref{sec:evaluation} reports the automatic discharge rate with
these helpers and the remaining explicit WD proofs. The measured
configuration includes this extension to the distributed automation.

\subsection{Individual obligations and reusable proofs}

Each import creates a collection of principal and WD obligations and
attempts automatic discharge. Users prove the remaining goals by name
or request the next pending obligation. The following excerpt is taken
from the completed case study. The ellipses omit other proved obligations;
the complete source contains no proof placeholders:
\par\noindent\begin{minipage}{\linewidth}
\begin{lstlisting}
import pog ZoneMonitor from "specs/zonemonitor"
-- ... other obligation commands ...
obligation Initialisation_0 of ZoneMonitor by
  intros
  exact ⟨Set.empty_subset _, by simp⟩
-- ... other obligation commands ...
qed ZoneMonitor
\end{lstlisting}
\end{minipage}\par
The \texttt{by} form accepts a tactic script, and \texttt{from} accepts a
proof term. The \texttt{of} clause selects an import; without it, the
latest import is used. Names are relative to the selected import,
including relative names of WD obligations.

A completed obligation is available to subsequent proofs in the same
import before final publication. A user can therefore establish a
function-image lemma or a common invariant consequence once and use it
in several proofs. Ordinary Lean definitions and lemmas can also be
introduced between obligation commands. Imports may be interleaved,
with each retaining its own progress and proof state. Diagnostics and
proof skeletons help locate residual goals and distinguish automatic
proofs, interactive proofs, admissions, and unsuccessful encodings.

Separate commands also allow Lean to retain unchanged command prefixes
during editing. Editing a late proof need not replay earlier proofs;
editing an earlier command causes the affected suffix to be elaborated
again. Tactics execute synchronously. These are properties of the proof
interface and its implementation, without implying automatic transport
of proofs after a change to the B specification.

\subsection{Completion and implementation}

An import retains a base and a working Lean environment. The latter
contains its generated declarations and completed proofs, which are
available while proving that import but remain unpublished in the
surrounding environment. The \texttt{qed} command checks that no obligation
is pending and no source goal failed to encode, then merges the checked
declarations into the current global environment. The merge rejects
name collisions, verifies replayed declarations against the kernel
environment, and restores metadata needed by generated definitions.
Nothing is published until that merge succeeds.

To avoid unnecessary replay, the implementation reuses a checked working
environment when a conservative comparison detects no external change.
New declarations or changed attributes can trigger a checked rebase.
Indexed obligation names, a pending-goal cursor, and maintained progress
counts support selection and reporting. These details implement the
same logical workflow for imports containing many obligations.

Automatic proof attempts receive individual heartbeat budgets, and
unsuccessful attempts leave their goals pending. Encoding failures are
recorded and block completion. Process-level failures such as exhaustion
of memory remain possible; successful import is consequently a separate
observation from complete discharge of the imported development.

\subsection{Trust boundary and proof completion}

Lean checks each proof against its translated statement and assumptions.
The connection to B additionally relies on Atelier~B's generation of the
obligations, the XML reader, extraction, and the encoding. A valid proof
of an incorrectly translated proposition would not establish the intended
source property. \barrel therefore provides kernel-checked proofs of
its generated statements, rather than a verified end-to-end translation.

The completion command follows Lean's convention for explicit admissions:
a proof containing \texttt{sorry} can be accepted with a warning. For a
completed case-study result, successful \texttt{qed} must therefore be
accompanied by an audit of the published declarations and their
dependencies, excluding admissions and unintended axioms. This separates
proof development in progress from the evidence reported in the paper.
